%bismi ALLAH ARRAHMAN ARRAHIM










































\documentclass[11pt,a4paper]{article}

% =========================================================
% PACKAGES
% =========================================================
\usepackage[utf8]{inputenc}
\usepackage[T1]{fontenc}
\usepackage[french]{babel}
\usepackage{lmodern}
\usepackage{microtype}

\usepackage{amsmath,amssymb,mathtools,bm}
\usepackage{booktabs,array,tabularx}
\usepackage{geometry}
\usepackage{xcolor}
\usepackage{graphicx}
\usepackage{float}
\usepackage{enumitem}
\usepackage{hyperref}

\usepackage{tikz}
\usetikzlibrary{
  arrows.meta,
  positioning,
  fit,
  calc,
  shapes.geometric,
  backgrounds
}

\geometry{
  left=2.2cm,
  right=2.2cm,
  top=2.2cm,
  bottom=2.2cm
}

\hypersetup{
  colorlinks=true,
  linkcolor=blue!55!black,
  urlcolor=blue!55!black,
  citecolor=blue!55!black
}

\setlength{\parindent}{0pt}
\setlength{\parskip}{5pt}
\setlist{nosep,leftmargin=1.5em}

% =========================================================
% COLORS AND TIKZ STYLES
% =========================================================
\definecolor{probeBlue}{RGB}{220,235,250}
\definecolor{matrixGreen}{RGB}{220,242,228}
\definecolor{rbfPurple}{RGB}{235,225,248}
\definecolor{coordOrange}{RGB}{250,232,210}
\definecolor{neutralGray}{RGB}{240,240,240}

\tikzset{
  block/.style={
    draw,
    rounded corners=2pt,
    thick,
    align=center,
    minimum height=0.85cm,
    text width=3.1cm,
    inner sep=5pt,
    font=\small
  },
  inputblock/.style={
    block,
    fill=neutralGray,
    text width=2.7cm
  },
  matrixblock/.style={
    block,
    fill=matrixGreen
  },
  rbfblock/.style={
    block,
    fill=rbfPurple
  },
  coordblock/.style={
    block,
    fill=coordOrange
  },
  arrow/.style={
    -{Latex[length=2mm]},
    thick
  },
  dashedarrow/.style={
    -{Latex[length=2mm]},
    thick,
    dashed
  }
}

% =========================================================
% DOCUMENT
% =========================================================
\title{
  \textbf{Formalisation des architectures EX7}\\
  \large Sondes RBF, sondes matricielles et coordination hybride
}

\author{Note technique de recherche}
\date{\today}

\begin{document}

\maketitle

\begin{abstract}
Cette note formalise quatre variantes d'une architecture composée
de sept sondes indépendantes et d'un mécanisme de coordination global :

\begin{enumerate}
  \item EX7-RBF : architecture de référence V003 ;
  \item EX7-Matrix : sondes locales matricielles ;
  \item EX7-Matrix-Vector : même architecture, avec calculs vectorisés ;
  \item EX7-Hybrid : sondes matricielles locales et coordinateur global RBF.
\end{enumerate}

L'objectif est de distinguer les effets de la représentation locale,
du mécanisme de coordination et de l'implémentation numérique.

Les architectures matricielles et hybrides présentées ici sont des
propositions expérimentales. Elles ne doivent pas être confondues
avec les résultats effectivement obtenus par EX7 V003.
\end{abstract}

\tableofcontents
\newpage

% =========================================================
\section{Notations générales}
% =========================================================

On considère un exemple d'entrée :

\[
  \bm{x}\in\mathbb{R}^{D}.
\]

Le système contient \(M\) sondes, avec ici :

\[
  M=7.
\]

Chaque sonde produit une représentation locale :

\[
  \bm{h}_i\in\mathbb{R}^{d_h},
  \qquad i\in\{1,\ldots,M\}.
\]

La concaténation des représentations est :

\[
  \bm{H}
  =
  \begin{bmatrix}
  \bm{h}_1\\
  \bm{h}_2\\
  \vdots\\
  \bm{h}_M
  \end{bmatrix}
  \in\mathbb{R}^{Md_h}.
\]

On note \(C\) le nombre de classes.

Pour EX7 V003, le problème expérimental possède :

\[
D=30,\qquad C=2,\qquad M=7.
\]

Les dimensions \(d_h\) des futures sondes matricielles seront
des hyperparamètres à définir et à comparer.

% =========================================================
\section{Architecture 1 : EX7-RBF de référence}
% =========================================================

\subsection{Organisation générale}

EX7 V003, également désigné
\textit{OuterCoord-Lite-7x1}, contient sept sondes
\texttt{AdaptiveRBF} indépendantes et un coordinateur global.

Chaque sonde reçoit le même vecteur d'entrée standardisé.
Les partitions d'entraînement sont disjointes, mais les sondes
ne reçoivent pas des sous-ensembles différents des dimensions.

Le schéma est le suivant :

\begin{figure}[H]
\centering
\begin{tikzpicture}[node distance=0.45cm]
  \node[inputblock] (x) {Entrée standardisée\\
    \(\bm{x}\in\mathbb{R}^{30}\)};

  \node[matrixblock,below=0.7cm of x,
    text width=10.5cm] (probes) {
    Sept sondes indépendantes\\[2pt]
    \(\operatorname{RBF}_1(\bm{x}),\ldots,
      \operatorname{RBF}_7(\bm{x})\)
  };

  \node[coordblock,below=0.7cm of probes] (rho) {
    Coordinateur global\\
    \(\bm{\rho}=\operatorname{softmax}(\bm{a})\)
  };

  \node[inputblock,below=0.7cm of rho] (out) {
    Logits et prédiction\\
    \(\bm{y}\in\mathbb{R}^{2}\)
  };

  \draw[arrow] (x) -- (probes);
  \draw[arrow] (probes) -- (rho);
  \draw[arrow] (rho) -- (out);
\end{tikzpicture}
\caption{Organisation conceptuelle d'EX7 V003.}
\end{figure}

\subsection{Fonction d'une sonde RBF}

Chaque sonde possède un encodeur appris :

\[
  \bm{z}_i
  =
  \tanh(W_i\bm{x}+\bm{b}_i),
\]

avec :

\[
  W_i\in\mathbb{R}^{d\times D},
  \qquad
  \bm{z}_i\in\mathbb{R}^{d}.
\]

Dans la configuration V003 :

\[
  d=4,\qquad K=8.
\]

La sonde possède \(K\) centres :

\[
  \bm{c}_{i,k}\in\mathbb{R}^{d}.
\]

La distance quadratique entre l'encodage et le centre \(k\) est :

\[
  \delta_{i,k}(\bm{x})
  =
  \left\|
    \bm{z}_i-\bm{c}_{i,k}
  \right\|_2^2.
\]

La largeur de chaque fonction est :

\[
  s_{i,k}=\exp(\ell_{i,k}).
\]

L'activation radiale est alors :

\[
  \phi_{i,k}(\bm{x})
  =
  \exp\left(
    -\frac{
      \delta_{i,k}(\bm{x})
    }{
      2s_{i,k}^{\,2}
    }
  \right).
\]

En regroupant les activations :

\[
  \bm{\phi}_i(\bm{x})
  \in\mathbb{R}^{K}.
\]

La sortie de la sonde est :

\[
  \bm{q}_i(\bm{x})
  =
  \bm{A}_i^\top\bm{\phi}_i(\bm{x})
  +\bm{c}_i^{\,out},
\]

où :

\[
  \bm{A}_i\in\mathbb{R}^{K\times C}.
\]

\subsection{Coordination globale}

Le coordinateur apprend un vecteur de poids :

\[
  \bm{\rho}
  =
  \operatorname{softmax}(\bm{a}),
  \qquad
  \bm{\rho}\in\mathbb{R}^{M}.
\]

Les sorties sont combinées par :

\[
  \boxed{
  \bm{y}(\bm{x})
  =
  \sum_{i=1}^{M}
  \rho_i\bm{q}_i(\bm{x})
  }
\]

Dans V003, \(\bm{\rho}\) est global et indépendant de l'entrée.
Il ne s'agit donc pas d'un mécanisme de sélection dynamique
des sondes.

\subsection{Propriété importante}

Les sondes sont entraînées sur des partitions stratifiées
disjointes des exemples. À l'inférence, les sept sondes
sont néanmoins toutes exécutées.

Il n'y a donc ni routage conditionnel, ni réduction du nombre
de sondes exécutées selon l'entrée.

% =========================================================
\section{Architecture 2 : EX7-Matrix}
% =========================================================

\subsection{Principe}

Cette variante conserve les sept unités locales indépendantes,
mais remplace les calculs RBF par des transformations matricielles.

Chaque sonde reçoit l'entrée complète :

\[
  \bm{x}\in\mathbb{R}^{D}.
\]

Elle calcule une représentation locale :

\[
  \boxed{
  \bm{h}_i
  =
  \sigma(W_i\bm{x}+\bm{b}_i)
  }
\]

avec :

\[
  W_i\in\mathbb{R}^{d_h\times D},
  \qquad
  \bm{b}_i\in\mathbb{R}^{d_h}.
\]

Une projection de sortie peut être ajoutée :

\[
  \bm{q}_i
  =
  V_i\bm{h}_i+\bm{c}_i,
\]

où :

\[
  V_i\in\mathbb{R}^{C\times d_h}.
\]

\subsection{Schéma}

\begin{figure}[H]
\centering
\begin{tikzpicture}[node distance=0.55cm]
  \node[inputblock] (x) {Entrée \(\bm{x}\)};

  \node[matrixblock,below left=1cm and 0.1cm of x] (p1) {
    Sonde 1\\
    \(W_1\bm{x}+\bm{b}_1\)\\
    \(\sigma(\cdot)\)
  };

  \node[matrixblock,below=1cm of x] (p2) {
    Sonde 2\\
    \(W_2\bm{x}+\bm{b}_2\)\\
    \(\sigma(\cdot)\)
  };

  \node[matrixblock,below right=1cm and 0.1cm of x] (p3) {
    Sonde 7\\
    \(W_7\bm{x}+\bm{b}_7\)\\
    \(\sigma(\cdot)\)
  };

  \node[coordblock,below=1.0cm of p2,
    text width=5.4cm] (coord) {
    Coordinateur matriciel\\
    \(\bm{y}=\sum_i\rho_i\bm{q}_i\)
  };

  \draw[arrow] (x) -- (p1);
  \draw[arrow] (x) -- (p2);
  \draw[arrow] (x) -- (p3);
  \draw[arrow] (p1) -- (coord);
  \draw[arrow] (p2) -- (coord);
  \draw[arrow] (p3) -- (coord);
\end{tikzpicture}
\caption{EX7-Matrix : calcul local matriciel et coordination.}
\end{figure}

Les sondes 3 à 6 ne sont pas dessinées pour conserver
la lisibilité du schéma.

\subsection{Coordination}

Une première possibilité consiste à conserver le coordinateur
pondéré d'EX7 :

\[
  \bm{\rho}
  =
  \operatorname{softmax}(\bm{a}).
\]

La sortie devient :

\[
  \boxed{
  \bm{y}
  =
  \sum_{i=1}^{M}
  \rho_i
  \left(
    V_i\sigma(W_i\bm{x}+\bm{b}_i)
    +\bm{c}_i
  \right)
  }
\]

Cette variante permet d'étudier l'organisation modulaire
sans conserver les fonctions radiales.

\subsection{Non-linéarité}

Si toutes les opérations sont uniquement affines,
leur composition reste affine.

En effet :

\[
  W_2(W_1\bm{x}+\bm{b}_1)+\bm{b}_2
  =
  W'\bm{x}+\bm{b}'.
\]

Une non-linéarité telle que ReLU, GELU ou tanh est donc
nécessaire pour obtenir une fonction globale non affine.

Pour un premier benchmark, ReLU constitue un choix simple :

\[
  \sigma(u)=\max(0,u).
\]

% =========================================================
\section{Architecture 3 : EX7-Matrix-Vector}
% =========================================================

\subsection{Différence avec EX7-Matrix}

EX7-Matrix-Vector ne change pas la fonction mathématique
des sondes. Il modifie la manière de calculer les sept branches.

On regroupe les matrices locales dans un tenseur :

\[
  \mathcal{W}
  \in
  \mathbb{R}^{M\times d_h\times D}.
\]

Pour un batch de \(B\) exemples :

\[
  X\in\mathbb{R}^{B\times D}.
\]

Le calcul des sondes peut être réalisé par une opération
batchée :

\[
  H
  =
  \sigma\left(
    \mathcal{W}\star X+\mathcal{B}
  \right),
\]

où \(\star\) désigne ici une multiplication matricielle
batchée, et :

\[
  H\in\mathbb{R}^{B\times M\times d_h}.
\]

\subsection{Schéma conceptuel}

\begin{figure}[H]
\centering
\begin{tikzpicture}[node distance=0.6cm]
  \node[inputblock] (x) {Batch d'entrée\\\(B\times D\)};

  \node[matrixblock,below=of x,
    text width=7cm] (tensor) {
    Tenseur des poids\\
    \(\mathcal{W}:M\times d_h\times D\)
  };

  \node[coordblock,below=of tensor,
    text width=7cm] (batch) {
    Multiplications matricielles batchées\\
    \(H:B\times M\times d_h\)
  };

  \node[inputblock,below=of batch] (out) {
    Fusion des sorties et prédiction
  };

  \draw[arrow] (x) -- (tensor);
  \draw[arrow] (tensor) -- (batch);
  \draw[arrow] (batch) -- (out);
\end{tikzpicture}
\caption{Calcul vectorisé des sondes.}
\end{figure}

La vectorisation peut réduire le coût des appels successifs
aux opérateurs numériques. Elle ne garantit toutefois pas
une accélération : le résultat dépend du matériel, de la
taille des batchs et de l'implémentation.

Il faut donc comparer les latences mesurées, et non déduire
la vitesse du seul nombre d'opérations.

% =========================================================
\section{Architecture 4 : EX7-Hybrid}
% =========================================================

\subsection{Hypothèse centrale}

Cette architecture correspond à la proposition hybride :

\begin{itemize}
  \item les sondes locales utilisent des transformations
  matricielles ;
  \item leurs représentations sont fusionnées ;
  \item le coordinateur global utilise une fonction RBF.
\end{itemize}

L'objectif est de séparer l'extraction locale des représentations
et la modélisation non linéaire de leurs configurations globales.

\subsection{Étape locale}

Pour chaque sonde :

\[
  \bm{h}_i
  =
  \sigma(W_i\bm{x}+\bm{b}_i).
\]

La fusion des représentations est :

\[
  \bm{H}
  =
  \operatorname{concat}
  (\bm{h}_1,\ldots,\bm{h}_M).
\]

Ainsi :

\[
  \bm{H}\in\mathbb{R}^{Md_h}.
\]

Une projection optionnelle peut être introduite :

\[
  \bm{u}=P\bm{H}+\bm{b}_P,
\]

avec :

\[
  P\in\mathbb{R}^{d_g\times Md_h}.
\]

Cette projection est facultative. Elle permet de contrôler
la dimension de l'espace dans lequel le coordinateur RBF
effectue ses calculs.

\subsection{Étape globale RBF}

Le coordinateur possède \(K_g\) centres :

\[
  \bm{\mu}_k\in\mathbb{R}^{d_g}.
\]

Les activations globales sont :

\[
  \psi_k(\bm{u})
  =
  \exp\left(
    -\frac{
      \|\bm{u}-\bm{\mu}_k\|_2^2
    }{
      2\tau_k^2
    }
  \right).
\]

On définit :

\[
  \bm{\psi}(\bm{u})
  =
  \begin{bmatrix}
  \psi_1(\bm{u})&
  \cdots&
  \psi_{K_g}(\bm{u})
  \end{bmatrix}^{\top}.
\]

La prédiction finale est :

\[
  \boxed{
  \bm{y}
  =
  A_g^\top\bm{\psi}(\bm{u})
  +\bm{b}_g
  }
\]

avec :

\[
  A_g\in\mathbb{R}^{K_g\times C}.
\]

\subsection{Schéma complet}

\begin{figure}[H]
\centering
\begin{tikzpicture}[node distance=0.65cm]
  \node[inputblock] (x) {
    Entrée\\
    \(\bm{x}\in\mathbb{R}^{D}\)
  };

  \node[matrixblock,below=of x,
    text width=8.2cm] (local) {
    Sept sondes matricielles\\
    \(\bm{h}_i=\sigma(W_i\bm{x}+\bm{b}_i)\)
  };

  \node[matrixblock,below=of local,
    text width=8.2cm] (fusion) {
    Fusion des représentations\\
    \(\bm{H}=\operatorname{concat}(\bm{h}_1,\ldots,\bm{h}_7)\)
  };

  \node[rbfblock,below=of fusion,
    text width=8.2cm] (rbf) {
    Coordinateur global RBF\\
    Distances, largeurs et activations radiales
  };

  \node[inputblock,below=of rbf] (out) {
    Projection finale\\
    \(\bm{y}\in\mathbb{R}^{C}\)
  };

  \draw[arrow] (x) -- (local);
  \draw[arrow] (local) -- (fusion);
  \draw[arrow] (fusion) -- (rbf);
  \draw[arrow] (rbf) -- (out);
\end{tikzpicture}
\caption{Architecture EX7-Hybrid proposée.}
\end{figure}

\subsection{Pourquoi cette séparation est intéressante}

Les transformations locales reposent sur des opérations
matricielles, tandis que le coordinateur RBF modélise
des configurations dans l'espace des représentations.

Cette organisation pourrait permettre :

\begin{enumerate}
  \item de conserver une partie du comportement non linéaire
  des RBF ;
  \item de simplifier les calculs des sondes ;
  \item de concentrer les calculs radiaux dans un seul module ;
  \item de comparer directement la coordination RBF
  et la coordination matricielle.
\end{enumerate}

Il ne s'agit cependant pas d'une garantie de gain de vitesse.
Le coordinateur global peut devenir un goulot d'étranglement,
en particulier si \(d_g\) ou \(K_g\) sont élevés.

% =========================================================
\section{Comparaison formelle des variantes}
% =========================================================

\begin{table}[H]
\centering
\small
\begin{tabularx}{\textwidth}{
  >{\raggedright\arraybackslash}p{2.4cm}
  X
  >{\centering\arraybackslash}p{2.2cm}
}
\toprule
Architecture & Calcul local & Coordination\\
\midrule
EX7-RBF
&
Encodeur, distances aux centres, exponentielles RBF
&
Pondération globale fixe\\

EX7-Matrix
&
Transformations matricielles et activation
&
Pondération globale\\

EX7-Matrix-Vector
&
Même fonction que Matrix, calculs batchés
&
Pondération globale\\

EX7-Hybrid
&
Transformations matricielles et activation
&
Coordinateur RBF global\\
\bottomrule
\end{tabularx}
\caption{Comparaison structurelle des quatre variantes.}
\end{table}

\subsection{Comptage des paramètres}

Pour EX7-Matrix, avec une couche cachée de dimension \(d_h\)
et une sortie de dimension \(C\), chaque sonde possède :

\[
  P_{\mathrm{sonde}}
  =
  d_hD+d_h+Cd_h+C.
\]

Pour \(M\) sondes et un coordinateur de \(M\) poids :

\[
  P_{\mathrm{Matrix}}
  =
  M P_{\mathrm{sonde}}+M.
\]

Pour EX7-Hybrid, si la projection globale est utilisée :

\[
  P_{\mathrm{Hybrid}}
  =
  M(d_hD+d_h)
  +d_g(Md_h+1)
  +K_gd_g
  +K_g
  +K_gC+C.
\]

Cette formule inclut :
\begin{itemize}
  \item les paramètres des sondes locales ;
  \item la projection globale ;
  \item les centres RBF ;
  \item les largeurs RBF ;
  \item la lecture finale et son biais.
\end{itemize}

Le comptage exact devra être adapté si certaines matrices
ou certains biais sont partagés.

Pour une comparaison équitable, il faudra compter également
le préprocesseur, les paramètres persistants et les éventuels
états nécessaires à l'inférence.

% =========================================================
\section{Relations entre les sondes et limites de la séparation}
% =========================================================

Dans les variantes précédentes, chaque sonde reçoit encore
l'entrée complète :

\[
  \bm{x}_i=\bm{x}.
\]

Cette propriété est volontaire pour la première comparaison :
elle permet de tester les mécanismes de calcul sans introduire
simultanément un découpage des dimensions.

Une évolution possible consisterait à distribuer les dimensions :

\[
  \bm{x}
  =
  \begin{bmatrix}
  \bm{x}_1\\
  \vdots\\
  \bm{x}_M
  \end{bmatrix}.
\]

Chaque sonde ne recevrait alors qu'un sous-espace :

\[
  \bm{h}_i
  =
  \sigma(W_i\bm{x}_i+\bm{b}_i).
\]

Cette variante peut réduire le coût local, mais un découpage
strict peut supprimer des interactions entre dimensions.

Il faudrait alors introduire une communication entre sondes,
des projections partagées ou des connexions croisées.

Le coordinateur RBF proposé précédemment ne doit pas être
considéré automatiquement comme une solution complète
à ce problème : il ne peut exploiter que les informations
qui lui sont transmises.

% =========================================================
\section{Protocole expérimental recommandé}
% =========================================================

\subsection{Comparaisons nécessaires}

Les quatre variantes doivent être entraînées et évaluées
dans des conditions identiques.

\begin{table}[H]
\centering
\small
\begin{tabularx}{\textwidth}{
  >{\raggedright\arraybackslash}p{3.0cm}
  X
}
\toprule
Mesure & Objectif\\
\midrule
Accuracy
&
Mesurer la proportion de prédictions correctes.\\

Log-loss
&
Évaluer la qualité probabiliste des prédictions.\\

Brier score
&
Mesurer l'erreur quadratique des probabilités.\\

Latence p50, p95, p99
&
Comparer le coût réel d'inférence.\\

Débit
&
Mesurer le nombre d'exemples traités par unité de temps.\\

Paramètres
&
Comparer la capacité paramétrique persistante.\\

Mémoire
&
Mesurer le coût mémoire du modèle et de son exécution.\\

Temps d'entraînement
&
Mesurer le coût de construction et d'optimisation.\\
\bottomrule
\end{tabularx}
\caption{Mesures à collecter.}
\end{table}

\subsection{Contrôle expérimental}

Les comparaisons devront conserver :

\begin{itemize}
  \item les mêmes partitions train et validation ;
  \item les mêmes graines aléatoires ;
  \item le même prétraitement ;
  \item les mêmes critères d'arrêt ;
  \item un budget de paramètres comparable ;
  \item le même matériel et les mêmes conditions de mesure.
\end{itemize}

Les résultats obtenus sur le jeu Wisconsin historique
ne constituent pas une validation externe indépendante.
Les différences entre graines sur un split fixe doivent
être interprétées comme des mesures descriptives.

Une validation sur des données indépendantes sera nécessaire
avant toute conclusion sur la généralisation.

% =========================================================
\section{Hypothèses scientifiques à tester}
% =========================================================

Les propositions peuvent être formulées sous forme
d'hypothèses réfutables.

\subsection*{H1 : efficacité des sondes matricielles}

À budget paramétrique comparable, EX7-Matrix pourrait
réduire la latence par rapport à EX7-RBF.

Cette hypothèse sera évaluée par les latences mesurées,
et non par le seul nombre de multiplications.

\subsection*{H2 : intérêt de la vectorisation}

EX7-Matrix-Vector pourrait réduire le surcoût d'exécution
des sept branches indépendantes.

Il faut mesurer séparément les batchs de taille 1 et les
batchs plus grands, car les gains peuvent varier selon
la taille des entrées.

\subsection*{H3 : intérêt de la coordination RBF globale}

EX7-Hybrid pourrait obtenir une meilleure modélisation
des interactions entre représentations locales qu'un
coordinateur linéaire.

Cette hypothèse exige une comparaison avec EX7-Matrix
à capacité paramétrique comparable.

\subsection*{H4 : spécialisation effective des sondes}

La séparation en sept sondes n'implique pas automatiquement
une spécialisation fonctionnelle.

Il faudra mesurer cette spécialisation à l'aide d'ablations,
de perturbations contrôlées des entrées et de comparaisons
des contributions individuelles des sondes.

% =========================================================
\section{Conclusion}
% =========================================================

La proposition EX7-Hybrid constitue une évolution cohérente
de l'architecture initiale.

Elle sépare deux responsabilités :

\begin{enumerate}
  \item l'extraction de représentations locales par des
  transformations matricielles ;
  \item la coordination globale par une fonction RBF.
\end{enumerate}

Cette séparation permet de tester si les fonctions RBF sont
plus utiles dans la coordination globale que dans chaque
sonde locale.

Il ne faut toutefois pas conclure à l'avance que cette
architecture sera plus rapide ou plus précise. Le coordinateur
RBF peut introduire un coût important, et la capacité du
système dépendra de la dimension des représentations,
du nombre de centres et de la qualité de la coordination.

L'expérience la plus informative consiste donc à comparer
EX7-RBF, EX7-Matrix, EX7-Matrix-Vector et EX7-Hybrid,
avec un protocole identique et des budgets paramétriques
contrôlés.

% =========================================================
\section*{Sources techniques}
% =========================================================

Les éléments décrivant l'implémentation de référence
proviennent de l'archive EX7 V003 :

\begin{enumerate}
  \item \texttt{reports/ARCHITECTURE\_V003.md} :
  organisation des sept sondes, shards d'entraînement,
  coordinateur global et protocole d'inférence.

  \item \texttt{inputs/reference\_v002/src/ex6\_common.py} :
  définition de la classe \texttt{AdaptiveRBF},
  encodeur latent, centres, largeurs et lecture RBF.

  \item \texttt{configs/PRE\_REGISTRATION\_V003.md} :
  hypothèses, protocole expérimental, comparateurs,
  mesures et limites méthodologiques.
\end{enumerate}

Les architectures Matrix, Matrix-Vector et Hybrid sont
des propositions formalisées dans cette note. Elles ne
sont pas des résultats expérimentaux déjà validés.

\end{document}
